\centerline{\bf \Huge Домашнее задание. Рыцари и лжецы}
\title{Домашнее задание. Рыцари и лжецы}

\begin{flushright}
        \scriptsize{Школа — это лишь тесный аквариум, \\а в мире есть океан и другие аквариумы. Выбери своё. \\ (c)Магическая битва. Наги Ёсино}
\end{flushright}

\problem{1} В лицее обучаются 88 восьмиклассников. Каждый из них может быть либо физматом, который всегда говорит правду, либо гуманитарием, который всегда лжет. Однажды все восьмиклассники выстроились по росту в шеренгу, и Начальник задал каждому вопрос: «Сколько гуманитариев ниже вас ростом?».\\
И получил следующие ответы:\\
1-ый(самый высокий): 77\\
2-ой: 77\\
3-ий: 77\\
$\dotsc$\\
82-ой: 77\\
83-ый: 77\\
84-ый: 3\\
85-ый: 3\\
86-ой: 2\\
87-ой: 2\\
88-ой(самый низкий): 1\\
Сколько физматов среди всех восьмиклассников лицея? 

\problem{2}За круглым столом сидят 10 учеников ЭлМШ. Каждый ученик может быть либо лососем, который всегда говорит правду, либо шиншиллой, которая всегда лжет, либо киви, которая может говорить правду, а может и лгать. Известно, что как минимум одна шиншилла сидит за столом. Начальник попросил каждого из учеников ответить на вопрос: «Кто сидит напротив вас?».\\
Первый ответил: «Напротив меня сидит шиншилла»;\\
А все остальные: «Напротив меня сидит лосось».\\
По полученым ответам Начальник смог определить все варианты количества шиншилл, лососей и киви, сидящих за столом.\\
А сможете ли вы?